\section*{Kapitel 3 --- Bussar, routing och tillförlitlighet}
\addcontentsline{toc}{section}{Kapitel 3 --- Bussar, routing och tillförlitlighet}

\solution{ex:3:1}{Routing}
\ans{a} I noden. En fysisk ledning levererar spänningsnivåer; den kan varken läsa adresser eller
veta att det finns frames. Varje nod måste därför själv jämföra \code{DST} med sin egen adress.
\ans{b} Ignorera den. Framen är inte felaktig --- den är adresserad till någon annan.
\ans{c} Därför att \code{Nack} betyder \emph{något gick fel}. På en buss med $n$ noder skulle varje
frame utlösa $n-2$ falska negativa kvittenser, och sändaren skulle inte kunna skilja dem från ett
verkligt fel.

\solution{ex:3:2}{Felmodellen}
\ans{a} \textbf{Tappad byte:} nej --- alla följande fält förskjuts, och \code{LEN} läses ur fel
byte. \textbf{Ändrad bit:} ja --- strukturen är intakt, bara innehållet är fel.
\textbf{Fördröjning:} ja --- ingenting i framen påverkas.
\ans{b} De två första fångas av checksumman, i \code{deserialize()}. Den tredje ska inte fångas
av något, eftersom ingenting är fel.
\ans{c} Efter en tappad byte eller en ändrad bit står parsern kvar med en förkastad frame och
måste nollställas; efter en fördröjning står den där den ska.

Fördröjningen ska alltså inte orsaka något problem. Gör den ändå det har parsern ett dolt
tidsberoende --- den antar något om \emph{när} bytes kommer --- och det är en bugg som inte visar
sig förrän systemet flyttas till långsammare hårdvara.

\solution{ex:3:3}{Förlorad kvittens}
\ans{a} A får ingen \code{Ack}, timern löper ut, och A skickar om framen med \textbf{samma}
\code{SEQ}.
\ans{b} B ser en frame den redan behandlat: samma avsändare, samma sekvensnummer.
\ans{c} B ska \textbf{skicka \code{Ack} igen} och \textbf{inte} köra applikationslogiken en andra
gång. Kommandot är redan utfört.
\ans{d} A hade fått timeout igen och skickat om igen, i en cykel som bara bryts av att
försöksräknaren tar slut. Det är därför den uteblivna kvittensen är det farliga felet, inte den
uteblivna framen.

\solution{ex:3:4}{Timeout-värdet}
\ans{a} Onödiga omsändningar och därmed fler dubbletter, plus ökad bussbelastning precis när
bussen redan är belastad.
\ans{b} Långsam felupptäckt: systemet står och väntar på ett svar som aldrig kommer, och
upplevs som hängt.
\ans{c} Värstafallstiden för en tur och retur: sändningstid, mottagarens behandlingstid och
kvittensens väg tillbaka. En nod känner sällan till den, eftersom den beror på hela systemets
belastning och på hur snabb den andra noden är --- vilket är just det som varierar.

\solution{ex:3:5}{Kedjan}
\textbf{Checksumma} upptäcker att en frame är trasig, men säger inget om frames som aldrig kom
fram. \textbf{ACK} bekräftar leverans, men hjälper inte när kvittensen själv uteblir.
\textbf{Timeout} upptäcker den uteblivna kvittensen, men gör inget åt den. \textbf{Omsändning}
åtgärdar den, och skapar därmed dubbletter. \textbf{Dubblettdetektering} städar upp efter
omsändningen, och skapar inget nytt problem --- det är därför kedjan slutar där.
